\section*{Проблематика оцифровки формул}
Вставка формулы осуществляется с использованием шрифтов, аналогичных Times New Roman. Размеры настроены автоматически:
\begin{equation}
	f(x) = \sum_{i=1}^{n} \frac{A_i}{B_{i^2}} + \alpha \beta \gamma
\end{equation}
Здесь индексы имеют размер 7 пт, а вложенные -- 5 пт, как того требуют правила ИрГУПС.

\section*{Пример вставки рисунка}
Рисунки должны быть черно-белыми или в градациях серого, 300 dpi. В журнал не принимаются цветные рисунки.
Если рисунок узкий, его можно вставить в одну колонку:

\begin{figure}[h]
	\centering
	% \includegraphics[width=\linewidth]{image_bw.jpg}
	\rule{\linewidth}{4cm} % Заглушка для рисунка
	\caption{Схема конвейера обработки}
	\label{fig:pipeline}
\end{figure}

Если рисунок широкий (до 150 мм), используйте окружение \texttt{figure*}, чтобы он растянулся на обе колонки:

\begin{figure*}[t]
	\centering
	% \includegraphics[width=150mm, height=100mm, keepaspectratio]{large_scheme_grayscale.tif}
	\rule{150mm}{5cm} % Заглушка для широкого рисунка
	\caption{Архитектура системы поиска и верификации (пример широкого рисунка)}
	\label{fig:architecture}
\end{figure*}

\section*{Оформление таблиц}
Подписи к таблицам оформляются шрифтом 10 пт, полужирным:

\begin{table}[h]
	\centering
	\caption{Сравнение метрик OCR и OCSR}
	\begin{tabular}{|p{3cm}|c|c|}
		\hline
		Метод & Precision & Recall \\
		\hline
		Tesseract & 0.85 & 0.82 \\
		DECIMER & 0.89 & 0.88 \\
		Наш конвейер & 0.94 & 0.91 \\
		\hline
	\end{tabular}
\end{table}

\section*{Заключение}
Выводы по работе. Общий объем статьи должен составлять 9–12 страниц формата А4.